\section{Constructing the finite system of a marked period chart}

% BEGIN L76F1
\paragraph{L76F1. Residues, vacuum lines and the multiplier.}
Fix the native symplectic period data, its labeling, and a level
$D$ with odd proper nonzero ideal $\mathfrak m$ coprime to
$\bar{\mathfrak m}$, residual frame $\beta_i$, group $J$, character
$\lambda$, and integral $w$ with $w+\bar w=1$. Write $L_*$ for the
period lattice of the chart, $S_L=\mathfrak m\bar{\mathfrak m}^{-1}L$,
and $A_{L,\Lambda}=\Lambda\otimes S_L$, in the matrix coordinates
of F74N2. Take $r\le n$, $n$ even, $M>0$ in $\mathfrak m$.
Use the original finite unitary Weil action and its Heisenberg
covariance. In L76F4--L76F5 we also use the original doubled
theta product and its Heisenberg intertwining, in the normalization
of O74E2.

At each prime of $\mathfrak m$ the residual frame is a basis:
the lattice is locally free of rank $r$, and its $r$ marked
generators span modulo the local ideal. Thus
$\Delta=D.\mathrm{grp}$ is exactly the set of rational unitary
$\gamma$ for which the matrix of $\gamma\beta$ reduces to some
$k_\gamma\in J$. Indeed such matrices are integral and invertible
on the local lattice, and reduction commutes with multiplication
and inversion, so the defining generators already form a subgroup.
The marked lattices $\mathcal L=D.\mathrm{lats}$ are the finitely
generated self-dual lattices agreeing with $L_*$ at the primes of
$\mathfrak m\bar{\mathfrak m}$. Local agreement is equivalent to
the common prime-to-ideal denominator in the source: apply
localization to the two finite quotient modules by the intersection.
Each $L$ spans the rational space by F74N1. Its transform under
$\gamma\in\Delta$ remains in $\mathcal L$, since unitarity transports
self-duality and duality at a prime of $\mathfrak m$ gives the
corresponding agreement at its conjugate prime.

For each $L$ choose one unit vacuum $v_L$ for
$A_{L,\mathcal O_E^n}$. Here is an existence and uniqueness
argument in the rational Schr\"odinger coordinates. Complete a
self-dual phase-space lattice $A$ at the finite places, let $U$
be its projection onto translations and $V$ its subgroup of pure
modulations. Self-duality gives $V=U^\perp$. The prescribed
semicharacter restricted to $V$ is a continuous character, hence
selects one coset $x_0+U$ of possible support. Indeed a character
on a compact open lattice factors through a finite quotient, where
the usual translation--modulation pairing identifies the characters
with the dual quotient; this also gives a rational representative
$x_0$. A value at $x_0$ determines the values on that coset by
lifting each translation in $U$ to $A$. Two lifts differ by $V$,
which acts with the required character there. The corrected
Heisenberg operators multiply exactly on $A$: the cross term of
$\epsilon_w(y)=\exp(\pi i\sum_i E_F(y_i,w y_i))$ cancels the
Heisenberg cocycle there by $w+\bar w=1$ and integral pairings.
The resulting nonzero function is locally constant (the phases
are trivial on a sufficiently small open sublattice) with compact
support. Restriction to rational points gives a Schwartz--Bruhat
vacuum, unique up to scalar. Normalize its norm to one.

Unitary $\gamma$ preserves both $E_F$ and $\epsilon_w$ and
carries $S_L$ onto $S_{\gamma L}$. Therefore there is a unique
unit $c(\gamma,L)$ with $w_G(\gamma)v_L=c(\gamma,L)v_{\gamma L}$.
Put $\mu(\gamma,L)=c(\gamma,L)/\lambda(k_\gamma)$.
Its norm is one, since $J$ is finite. If the chart supplies
vacua and a multiplier satisfying the native transport law, this
formula recovers precisely that multiplier. For every allowed
$\Lambda$, let $u_L$ be the normalized orthogonal projection of
$v_L$ to the $A_{L,\Lambda}$-vacuum line. Its projection is
nonzero by F74N2. Unitary transport commutes with the projector
and the real normalization, so it takes $u_L$ to
$c(\gamma,L)u_{\gamma L}$. In the native finite sum
\[
   \Phi_L=\sum_{k\in J}\lambda(k)T_{y_k}u_L
\]
use the representatives of F74N2. Transport takes $y_k$ to a
representative of $k_\gamma k$. Reindexing contributes
$\lambda(k_\gamma)^{-1}$, which gives
$w_G(\gamma)\Phi_L=\mu(\gamma,L)\Phi_{\gamma L}$. Each $\Phi_L$
is Schwartz--Bruhat by its finite expression.
% END L76F1

% BEGIN L76F2
\paragraph{L76F2. Orthogonality outside the level subgroup.}
For rational unitary $\gamma$ suppose a summand of
$\langle\Phi_L,w_G(\gamma)\Phi_{L'}\rangle$ is nonzero, with
representatives $y_k,y'_{k'}$. Write $A=A_{L,\Lambda}$ and
$A'=A_{L',\Lambda}$, using $\gamma$ on each row. The transported
vacuum is a vacuum for $\gamma A'$. Translating both summands by
any $z\in A\cap\gamma A'$ cancels their common vacuum character.
The commutator character of the difference of their shifts must
then be trivial. Duality as in F74N2 gives
\[
 y_k-\gamma y'_{k'}\in(A\cap\gamma A')^\vee=A+\gamma A'.
\]
Consequently $y_k-a=\gamma(y'_{k'}-a')$ for some $a\in A$,
$a'\in A'$. At every prime of $\mathfrak m$ the rows of $a,a'$
are in $\mathfrak m L_*$ locally. The first $r$ rows on each
side before applying $\gamma$ are local bases with residue
matrices $k,k'$, respectively. Thus $\gamma$ preserves the
reference lattice at those primes, with reduction $k(k')^{-1}\in J$.
By L76F1 this puts $\gamma$ in $\Delta$. Summing the finitely
many inner products proves the support field for every
$\gamma\notin\Delta$, including matrices with denominators at
the marked primes.
% END L76F2

% BEGIN L76F3
\paragraph{L76F3. An annihilator in every rational isotropic direction.}
Fix $0\ne\xi\in E^r$ with $h(\xi,\xi)=0$ and an allowed $L,\Lambda$.
On the additive group $B=\{v:v\otimes\xi\in A_{L,\Lambda}\}$,
the operator $T_{v\otimes\xi}$ acts on $T_yu_L$ with character
\[
  \chi_y(v)=\exp\!\left(-2\pi i\sum_j E_F(v_j\xi,y_j)\right).
\]
The semicharacter of $v\otimes\xi$ is one by isotropy, and
commuting its displacement past $T_y$ gives the displayed phase.

For every representative with nonzero coefficient this character
is nontrivial. To prove it, suppose the row pairing in the
exponent is integral for all $v\in B$. Choose a positive integer
$t$ prime to $M$ with $t\mathcal O_E^n\subset\Lambda$.
For every $z\in E\xi\cap S_L$ and every row $j$, testing the
one-row displacement $t e_j\otimes z$ gives $E_F(tz,y_j)\in
\mathbb Z$. The same holds against the integral span of the
first $r$ rows: a scalar on the second argument moves to its
conjugate on the first, preserving $E\xi\cap S_L$. Their span
contains $sL_*$ for some $s\in\mathcal O_E$ prime to
$\mathfrak m$, since they are bases at its primes. Indeed the
annihilator of $(L_*+\sum\mathcal O_E y_j)/\sum\mathcal O_E y_j$
is coprime to $\mathfrak m$, so $s$ can be chosen congruent to one.
Self-duality now gives $t\bar s z\in L_*$. Choose $b$ prime to
$\mathfrak m\bar{\mathfrak m}$ with $b L_*\subset L$. Thus
$u=b t\bar s$ carries $E\xi\cap S_L$ into $L$. At any prime
$P\mid\bar{\mathfrak m}$ it is a unit: $t$ is prime to $M$ and
the integer $M$ belongs to both ideals. Write
$e=\operatorname{ord}_P(\bar{\mathfrak m})>0$. Locally
$S_{L,P}=\pi_P^{-e}L_P$, so intersecting with the line gives the
claimed inclusion of $\pi_P^{-e}(E\xi\cap L_P)$ into
$E\xi\cap L_P$. This is impossible for a rank-one lattice over
a discrete valuation ring.

For each of the finitely many $y$ choose $v_y\in B$ with
$\chi_y(v_y)\ne1$. In the group algebra of $B$ take the product
of $[v_y]-\chi_y(v_y)[0]$. Its coefficient sum is
$\prod_y(1-\chi_y(v_y))\ne0$. Its action annihilates every
summand, since each has its own vanishing factor. These operators
compose as the additive group because the row symplectic form
vanishes on the isotropic line. Expanding the product gives exactly
the finite coefficients required by the cusp field of the system.
% END L76F3

% BEGIN L76F4
\paragraph{L76F4. The doubled vacuum has one phase for every definite lattice.}
Put $S_1=S_L$, $S_2=S_{L'}$, $K=S_1+S_2$, and
$H=(\kappa_0/2)G$, the native invertible skew-Hermitian matrix.
Write $\mathcal D$ for doubling. For the normalized projections
$u_1,u_2$ of the two fixed standard-lattice vacua $v_1,v_2$,
there are $C=C^*$ over $E$ and $t$ independent of $\Lambda$ such that
\[
 \mathcal D(u_1,u_2)(x)=
 t\,\mathbf1_{\Lambda\otimes K}(x)
 \exp\!\big(\pi i\operatorname{Tr}_{E/\mathbb Q}
                   \operatorname{tr}(xCx^*)\big).
\]

First compute the doubled phase-space lattice. The change of
coordinates is $(y_1,y_2)\mapsto(d,b)=(y_1-y_2,(y_1+y_2)H/2)$.
For this paragraph use row vectors. The projection onto $d$ is
$K$ and its vertical kernel is $(S_1\cap S_2)H=K^\perp$ for
the trace of the ordinary Hermitian dot product: self-duality of
each $S_i$ identifies its integral symplectic dual with itself.
The surjection $S_1\oplus S_2\to K$ splits because $K$ is a
finitely generated torsion-free module over a Dedekind domain,
hence projective. Extend the section rationally, writing its
components as $dA,d(A-I)$. Its graph in the doubled coordinates
has matrix $C_0=(A-I/2)H$. Set $C=(1-w)C_0+wC_0^*$.
Then $C=C^*$, and
\[
 C_0^*-C_0=(A-I)H(A-I)^*-AHA^* .
\]
When paired with $d,d'\in K$ the trace of
$d\,w(C_0^*-C_0)d'^*$ is integral. The displayed difference
pairs elements of $S_2$ with $S_2$, and of $S_1$ with $S_1$,
after multiplication by the integral $w$; both are integral by
self-duality. Thus this correction is vertical and the same
lattice has the section $d\mapsto(d,dC)$. Tensoring this section
with $\Lambda$ gives lifts for every $x\in\Lambda\otimes K$.
The pure modulations, from $y_1=y_2$, force the doubled vacuum
to vanish outside this tensor lattice. Indeed the character
there must pair integrally with the intersection of the two
self-dual tensor lattices, whose dual is their sum.

On the section, the two vacuum semicharacters cancel. In one row,
$dC=(y_1+y_2)H/2$ and Hermiticity makes $dCd^*$ real.
Subtracting its conjugate yields $y_1Hy_1^*=y_2Hy_2^*$ by
$H^*=-H$. The same therefore holds for the traced forms
$E_F(y_i,w y_i)$. Applying the doubled Heisenberg transformation
for the section displacement and evaluating at $x=d$ gives the
value at zero times the displayed quadratic exponential. The same
calculation works for arbitrary unit vacua at the two tensor
lattices and in particular for the standard pair.

It remains to identify the value at zero across $\Lambda$,
including its argument. The finite Heisenberg laws imply for
Schwartz--Bruhat tests, with inner product linear in the first slot,
\[
 \langle\mathcal D(f,g),\mathcal D(f',g')\rangle
    = b\langle f,f'\rangle\overline{\langle g,g'\rangle},
 \qquad b>0 .
\]
For fixed $g,g'$ the left side is invariant under simultaneous
Heisenberg displacement of $f,f'$ by the doubled intertwining and
the unitarity of Hermitian translations and modulations. The coset
indicator argument of O74E2 makes it proportional to $\langle
f,f'\rangle$. Apply the same argument to the second pair in the
conjugate order. The proportionality constant is positive, since
the double of the standard indicator with itself is nonzero by
O74E2. Applying this identity to $(v_1,v_2)$ and $(u_1,u_2)$
gives a positive real number: $\langle v_i,u_i\rangle$ is the
positive norm of its nonzero orthogonal projection. On the other
hand the already established formula for the two functions gives
this inner product as $t_0\bar t_\Lambda$ times the positive
volume of the intersection of the supports. Hence their phases
agree. By O74E2, their absolute values are
$|a|\,|\langle v_1,v_2\rangle|$ and
$|a|\,|\langle u_1,u_2\rangle|$. Both overlaps equal the
inverse square root of the same index
$[S_1:S_1\cap S_2]^n$ by F74N2. This proves the assertion with
$t=a\langle v_1,v_2\rangle$, including its phase.
% END L76F4

% BEGIN L76F5
\paragraph{L76F5. Admission of the actual doubled finite sum.}
Keep one pair $L,L'$ fixed. Put
$K_0=2\bar{\mathfrak m}(L\cap L')$, $B_0=H/2$ and
$A_0=\mathcal O_E^r$. For each $\Lambda$ take common representatives
$c_k\in\Lambda\otimes K_0$ with first $r$ rows reducing to $k$
and remaining rows zero modulo $\mathfrak m$. By the
representative independence in F74N2 we may use any such choices.
By sesquilinearity and the same change of Heisenberg coordinates,
the doubled test is a finite sum, with coefficients
$\lambda(k)\overline{\lambda(k')}$, of
\[
 T_H\big(c_k-c_{k'},(c_k+c_{k'})B_0\big)
 \left[t\,\mathbf1_{\Lambda\otimes K}(x)
     \exp\!\big(\pi i\operatorname{Tr}_{E/\mathbb Q}
                         \operatorname{tr}(xCx^*)\big)\right].
\]
All constants and coefficient lattices here precede the choice
of $\Lambda$. Write $XY$ for the image of a row lattice $X$ by
right multiplication by $Y$. Clear denominators with positive
integers $N,N'$ so that $N$ carries
$K,K_0,K_0B_0,KC,KC^*,A_0 C$ into $A_0$ and $N'A_0\subset K$.
Fullness of $K$ and finite generation supply these integers.
Then $R=2NN'$ works in the definition of admission.

Indeed a nonzero summand has $x-d\in\Lambda\otimes K$ with
$d=c_k-c_{k'}\in\Lambda\otimes K_0$, whence $Rx$ has columns
in $\Lambda$. For a matrix $y$ with columns in $\Lambda$,
$Ry\in\Lambda\otimes K$ so translation by $Ry$ preserves the
indicator. In the exponent of the quadratic at $z=x-d$ it adds
the traces of $RyCz^*$, $RzCy^*$ and $R^2yCy^*$. The denominator
conditions on $KC^*,KC,A_0C$ and integral Hermitian pairings on
$\Lambda$ put each trace in $2\mathbb Z$. The modulation changes
by the trace of $2R(c_k+c_{k'})B_0y^*$, again even integral.
These are the exact period and support requirements. For a unitary
row move congruent to the identity at $M$ taking $\Lambda$ to
$\Lambda'$, the transported $hc_k$ are permissible representatives.
Indeed locally at $\mathfrak m$, $h-I$ has entries in
$M\mathcal O_{E,P}$ and all the rows of $c_k$ are integral in
the reference lattice. Every trace in the formula is invariant by
$h^*h=1$ and the support lattice is transported to the new one.
This proves equivariance.

Finally assume the two lattices agree at $RM$, and $Rx$ has
columns in both. Choose $s>0$ prime to $RM$ with
$s\Lambda\subset\Lambda'$ and $s\Lambda'\subset\Lambda$.
Replacing every $c_k$ for $\Lambda$ by
$s^{\varphi(M)}c_k$ gives permissible representatives for both:
the factor is one modulo $M$ and is a positive power of $s$.
With these common choices only the support indicator can differ.
For $z=x-d\in\Lambda\otimes K$, one has
$sz\in\Lambda'\otimes K$; also
$N'Rz\in\Lambda'\otimes K$ since $Rx,Rd$ have columns in
$\Lambda'$. The integers $s,N'R$ are coprime because $N'\mid R$.
B\'ezout gives $z\in\Lambda'\otimes K$, and interchanging the
lattices gives the converse. This proves locality on precisely
the common allowed support. The single $R$ for this pair proves
the native admission field.
% END L76F5

% BEGIN L76F6
\paragraph{L76F6. The finite input to spanning on a supplied chart.}
Take even $n\ge4[E:\mathbb Q]r+4$. Then the data of L76F1
give the source system
$(M,\Delta,\mathcal L,D.\mathrm{phi},\mu)$ on the fixed chart,
under the finite transformation laws there specified. The group,
lattice, stability, Schwartz--Bruhat and multiplier fields are
L76F1. F74N2--F74N3 give diagonal $|J|>0$, majorant
$[L:L\cap L']^{-n/2}$ and uniform off-diagonal sum at most
$0.65$. L76F2, L76F3 and L76F5 give its remaining three fields.
Thus the finite system required in I75K3 can be constructed from
the given marked level and compatible period data. In applications
the automorphy law for the form must use the multiplier determined
by the chosen vacuum phases in L76F1. Construction of a common
principal chart and comparison of its classical theta functions
with the geometric short sections require the original geometric
data and comparison laws.
% END L76F6
